\begin{lemma}[A priori bound on $\psi$]
  Let $E$ be a metric space, let $m \in \mathbb{N}$ with $m \ge 1$, let $Q \subset E$ be a finite
  set, let $\epsilon, C \in \mathbb{R}$ with $\epsilon \ge 0$ and $C \ge 0$, and let
  $\gamma \in \Theta(E)$ (\noderef{Curve}). Then
  \[
    0 \;\le\; \psi_{m,Q,\epsilon,C}(\gamma)
      \;\le\; 2C^2\length(\gamma) + \frac{2C^2\length(\gamma)^2}{m}
    ,
  \]
  where $\psi_{m,Q,\epsilon,C}$ is as in \noderef{psi}. In particular the upper bound depends only
  on $m$, $C$ and $\length(\gamma)$: neither $Q$ nor $\epsilon$ nor any further feature of $\gamma$
  enters it.

  \paragraph{Why this node is needed.} The strong-law assembly of Theorem 4.4 (paper.tex
  1349--1360) applies the Strong Law of Large Numbers to the countable family of functions
  $\gamma \mapsto l^{-1}\psi_{m\cdot l,Q,\epsilon,C}(\gamma)$ on $\Theta(E)$, one for each
  admissible $(m,Q,\epsilon,C)$ and each stage $l$, and the Strong Law requires each member of that
  family to be $\overline\eta_l$-integrable. Measurability is \noderef{PsiMeasurable}; integrability
  against the \emph{finite} measure $\overline\eta_l$ needs exactly a uniform bound of the present
  shape, since on the $\overline\eta_l$-a.e.\ event $\length(\gamma)=l$ the displayed bound at
  $m\cdot l$ reads $2C^2l + 2C^2l^2/(m l) = 2C^2 l + 2C^2 l/m$, so that
  $l^{-1}\psi_{m\cdot l,Q,\epsilon,C}$ is a.e.\ bounded by the $l$-independent constant
  $2C^2 + 2C^2/m$. Stating the bound at a general $m$ and a general $\length(\gamma)$, rather than
  specialising to $(m\cdot l, l)$, is what makes it reusable at every stage $l$ simultaneously.
\end{lemma}
\begin{proof}
  Recall \noderef{psi}'s defining formula: writing $L := \length(\gamma)$ and, for
  $i \in \{1,\dots,m\}$,
  \[
    u_i := \min\Bigl\{\tfrac{2CL}{m},\;
      2\epsilon + 2C\bigl(\mathrm{infDist}(\gamma(\tfrac{i}{m}L),Q)
        + \mathrm{infDist}(\gamma(\tfrac{i-1}{m}L),Q)\bigr)\Bigr\}
    ,
  \]
  we have $\psi_{m,Q,\epsilon,C}(\gamma) = C\sum_{i=1}^m u_i + \frac{2C^2L^2}{m}$.

  \emph{Preliminaries.} Since $m \ge 1$ we have $m > 0$ as a real number, and
  $L = \length(\gamma) \ge 0$ by \noderef{Curve} (field \emph{len\_nonneg}). Since $C \ge 0$ and
  $L \ge 0$, the first argument of each minimum satisfies $\tfrac{2CL}{m} \ge 0$. Also
  $\tfrac{2C^2L^2}{m} \ge 0$, being a quotient of nonnegative reals by a positive real.

  \emph{Lower bound.} Fix $i \in \{1,\dots,m\}$. The first argument of $u_i$'s minimum is
  $\ge 0$, as just noted. For the second argument, $\mathrm{infDist}(x,Q) \ge 0$ for every $x$ (an
  infimum of distances, each $\ge 0$), and $\epsilon \ge 0$, $C \ge 0$, so
  $2\epsilon + 2C(\mathrm{infDist}(\gamma(\tfrac{i}{m}L),Q) +
  \mathrm{infDist}(\gamma(\tfrac{i-1}{m}L),Q)) \ge 0$ as well. A minimum of two nonnegative reals is
  nonnegative, so $u_i \ge 0$. Hence $\sum_{i=1}^m u_i \ge 0$, and since $C \ge 0$ also
  $C\sum_{i=1}^m u_i \ge 0$. Adding the nonnegative term $\tfrac{2C^2L^2}{m}$ gives
  $\psi_{m,Q,\epsilon,C}(\gamma) \ge 0$.

  \emph{Upper bound.} Fix $i \in \{1,\dots,m\}$. A minimum is at most its first argument, so
  $u_i \le \tfrac{2CL}{m}$. Summing this inequality over the $m$ indices
  $i \in \{1,\dots,m\}$ (the index set $\{1,\dots,m\}$ has exactly $m$ elements),
  \[
    \sum_{i=1}^m u_i \;\le\; m \cdot \frac{2CL}{m} \;=\; 2CL
    ,
  \]
  the last equality because $m > 0$. Multiplying by $C \ge 0$ preserves the inequality:
  $C\sum_{i=1}^m u_i \le 2C^2L$. Adding $\tfrac{2C^2L^2}{m}$ to both sides,
  \[
    \psi_{m,Q,\epsilon,C}(\gamma) = C\sum_{i=1}^m u_i + \frac{2C^2L^2}{m}
      \;\le\; 2C^2L + \frac{2C^2L^2}{m}
    ,
  \]
  which is the asserted upper bound.
\end{proof}
